\section{本章小结}
\label{sec:opt-summary}

本章把“在允许的选择里找最好的”拆成了几个分别回答的问题。写一个优化问题要同时给出决策变量、目标函数和可行集。最优值是下确界，最优解是取得它的可行点，两者可以分离：问题可能不可行、无下界，或最优值有限却未被达到。局部最优只和附近的可行点比较，全局最优和所有可行点比较。约束指标函数把约束并入目标而不改变问题，普通指示函数和有限惩罚则不能这样用。

存在性靠最小化序列的三关：留得住、极限合法、比较传得到。紧性通过前两关，下半连续通过第三关，无界可行集用紧下水平集或强制性处理。解的性质靠线段比较：凸性使局部最优成为全局最优，严格凸使解至多一个，强凸使目标差控制到解的距离。

判别依赖一个点的导数信息。一阶和二阶条件在内点上区分驻点、局部最优和鞍点，它们本身是局部的；只有目标凸时，驻点才是全局最优。边界上的最优点梯度可以不为零，因为允许的方向变少了。

两个贯穿的例子展示了这些问题怎样得到回答。资源分配的可行集是有界多面体，解存在；把两条资源约束相加，同时得到下界和达到它的点$(2,2)$。二次拟合的$\gamma>0$让目标强制，又让Hessian正定，解存在且唯一，目标差可以换成距离；$\gamma=0$且$\bm{A}$秩亏时，目标不再强制，解仍存在，却成了一条直线。

表~\ref{tab:opt-summary}汇总各结论所需的条件，表中均为充分条件，反例列说明去掉条件或反向推断时的后果。

\begin{table*}[!htbp]
\centering
\small
\begin{tabularx}{\linewidth}{@{}>{\RaggedRight\arraybackslash}p{24mm}>{\RaggedRight\arraybackslash}X>{\RaggedRight\arraybackslash}X>{\RaggedRight\arraybackslash}X@{}}
\toprule
问题 & 充分条件 & 结论 & 反例或边界 \\
\midrule
最优解是否存在 & 可行集非空紧，目标下半连续且有有限值；或某个下水平集非空紧 & 最优值有限且被达到 & $(0,1]$上的$x$；$[1,\infty)$上的$1/x$；在$0$处跳高的函数 \\
局部是否即全局 & 可行集凸，目标凸 & 局部最优解都是全局最优解，解集凸 & 例子中的四次函数；$[-3,-1]\cup[1,3]$上的$(x-2)^2$ \\
解是否唯一 & 可行集凸，目标严格凸，解存在 & 解唯一 & $e^x$严格凸而无最小点；$\lvert x\rvert$与资源分配有唯一解而不严格凸 \\
目标差能否换成距离 & 可行集凸，目标$\mu$强凸，解存在 & $f(\bm{x})-f^\star\geq\frac{\mu}{2}\lVert\bm{x}-\bm{x}^\star\rVert_2^2$ & 非凸可行集$\{\lvert x\rvert\geq1\}$上的$x^2$ \\
内点能否用梯度判别 & 可微；开凸集上目标凸 & 局部最优必有$\nabla f=\bm{0}$；凸时$\nabla f=\bm{0}$即全局最优 & 鞍点；边界最优点的导数可不为零 \\
驻点是否局部最优 & 二阶连续可微，Hessian正定 & 严格局部最优解 & $x_1^2+x_2^4$（正定非必要）；$x_1^2+x_2^3$（半正定不充分）\\
小梯度是否表示小误差 & $\mu\bm{I}\preceq\nabla^2f\preceq L\bm{I}$，无约束 & $e_f\leq e_g^2/(2\mu)$，$e_d\leq e_g/\mu$ & 平坦目标$0.005x^2$；$\gamma=0.01$的拟合实例 \\
\bottomrule
\end{tabularx}
\caption{本章主要问题、充分条件与边界。}
\label{tab:opt-summary}
\end{table*}

算法的输出和问题的解是两回事。输出$\hat{\bm{x}}$由目标差、解距离、梯度大小、可行性残差和最优性残差分别衡量，它们回答不同的问题，只有在曲率上下界这样的条件下才能换算，系数随$1/\mu$放大。

本章没有构造任何算法。从下一章起，线性规划把这些语言用在最具体的模型上：可行集是多面体，解在顶点中寻找，“给约束配权重再相加”发展成可以核对的最优性证书。
